\title{Controlling Text-to-Image Diffusion by Orthogonal Finetuning}

\begin{abstract}
Recent advancements in large text-to-image diffusion models have enabled remarkable synthesis of realistic images from textual descriptions. A central challenge now lies in determining how best to guide and modify these robust models for a variety of downstream applications. In response, we propose a systematic approach to finetuning, termed Orthogonal Finetuning~(OFT), designed specifically for adapting text-to-image diffusion models to specific tasks. Distinct from prior approaches, OFT can be theoretically shown to maintain hyperspherical energy, a metric representing the pairwise correlations among neurons situated on a unit hypersphere. Our analysis reveals that retaining this property is essential for upholding the model’s semantic image generation capabilities. To further enhance the robustness of finetuning, we introduce Constrained Orthogonal Finetuning (COFT), which incorporates an extra constraint on the hypersphere's radius. We investigate two key applications of finetuning in text-to-image generation: subject-driven generation—aiming to produce images of particular subjects based on a handful of input images and a prompt; and controllable generation—enabling the model to process auxiliary control inputs. Through empirical studies, we demonstrate that the OFT methodology not only surpasses existing techniques in terms of output fidelity but also achieves faster convergence.
\end{abstract}

\section{Introduction}

Recent advancements in text-to-image diffusion models~\cite{saharia2022photorealistic,ramesh2022hierarchical,rombach2022high} have led to remarkable results in generating high-quality images guided by textual prompts. However, even with these successes, textual instructions often remain vague or lack the necessary specificity for detailed and precise manipulation of generated content. In response to these limitations, this work focuses on two specific categories of text-to-image generation tasks:

\item \textbf{Subject-driven generation}~\cite{ruiz2023dreambooth}: This task involves synthesizing images of a specific subject, for which only a limited set of images is available, and placing the subject into new contexts by utilizing a text prompt as guidance.
\item \textbf{Controllable generation}~\cite{zhang2023adding,mou2023t2i}: Here, image synthesis is guided not only by a textual description but also by an extra control input (\emph{e.g.}, canny edges or segmentation maps), requiring the generated output to conform to both sources of guidance.
\end{itemize}

At their core, both tasks concern the challenge of effectively adapting text-to-image diffusion models through finetuning, while maintaining the strong generative performance achieved during pretraining. The key requirements for a robust finetuning approach can be summarized as follows: (1) \emph{training efficiency}, characterized by a reduction in trainable parameters and training epochs; and (2) \emph{generalizability preservation}, which refers to retaining the capacity for high-fidelity and diverse generation. Standard finetuning techniques approach this by either incrementally modifying the neuron weights with a modest learning rate (\emph{e.g.}, \cite{ruiz2023dreambooth}), or by incorporating an auxiliary module with re-parameterized neuron weights (\emph{e.g.}, \cite{hulora2022,zhang2023adding}). Although straightforward, these strategies do not inherently guarantee that the pretrained generative abilities will be conserved. Furthermore, the field currently lacks a systematic framework for formulating optimal finetuning methodologies or selecting appropriate hyperparameters, such as the duration of training. A significant obstacle arises from the absence of a quantitative criterion to evaluate the extent to which generative abilities are preserved post-finetuning. Conventional approaches often make the tacit assumption that a lesser Euclidean distance between the finetuned and pretrained model weights equates to better retention of the original model's capabilities. This notion underlies the widespread use of minimal learning rates in finetuning regimes. Nevertheless, we contend that relying solely on Euclidean distance as a metric is insufficient to assess semantic preservation comprehensively. Instead, introducing a more structural metric to evaluate differences between the finetuned and pretrained models could lead to improved retention of pretrained capabilities and enhance the overall stability of the finetuning process.

Drawing on empirical findings that hyperspherical similarity effectively represents semantic relationships~\cite{liu2017deep,liu2018decoupled,chen2020angular}, we leverage the concept of hyperspherical energy~\cite{liu2018learning} to describe the pairwise interactions among neurons. Hyperspherical energy aggregates the hyperspherical similarity (such as cosine similarity) across all neuron pairs within a given layer, providing a quantitative measure of how uniformly neurons are distributed on the unit hypersphere~\cite{liu2021learning}. Our central assumption is that an optimally finetuned model should exhibit minimal alteration in hyperspherical energy relative to the pretrained baseline. Although one could attempt to enforce hyperspherical energy invariance during finetuning via a regularization term, such an approach does not guarantee the reduction of hyperspherical energy discrepancy. Instead, we exploit a notable invariance property: pairwise hyperspherical similarity is mathematically conserved under uniform orthogonal transformations applied to all neurons. Building on this property, we introduce Orthogonal Finetuning~(OFT), a methodology aimed at adapting large-scale text-to-image diffusion models for downstream tasks while maintaining their hyperspherical energy. The approach centers on learning a single orthogonal transformation per layer to preserve the pairwise angular relationships among neurons. Conceptually, OFT corresponds to selecting a different canonical basis for the neurons in each layer. Furthermore, recognizing that reducing the Euclidean distance between the pretrained and finetuned models generally favors retention of pretrained capabilities, we advance Constrained Orthogonal Finetuning~(COFT), an extension of OFT that confines the finetuned model to a hypersphere of constant radius centered at the pretrained configuration.

The reasoning behind the effectiveness of orthogonal transformations in finetuning neurons is partially analogous to principles found in the 2D Fourier transform, wherein an image can be represented by its magnitude and phase spectra. The phase spectrum encapsulates the angular relationship between the input and its basis functions, and is chiefly responsible for retaining the semantic content of the image. Empirically, reconstructing an image with its original phase spectrum combined with an arbitrary magnitude spectrum tends to preserve the semantic integrity of the original image. This insight indicates that the orientation of neurons plays a critical role in modifying the semantic attributes of generated images, which aligns with the objectives of subject-driven and controllable image generation. Nevertheless, introducing unbounded flexibility in altering neuron directions can compromise the generative capability established during pretraining. To regulate this flexibility, we advocate for the preservation of inter-neuronal angles, grounded in the hypothesis that these angular relationships encapsulate essential knowledge within neural networks. Following this rationale, we propose to implement layer-wise orthogonal transformations shared among all neurons in each layer, thereby maintaining constant hyperspherical energy.

Our approach is also motivated by orthogonal over-parameterized training~\cite{liu2021orthogonal}, which introduces a methodology for training classification neural networks from the ground up through orthogonal transformations applied to randomly initialized networks. This direction leverages the fact that random initialization results in neural networks with demonstrably low hyperspherical energy in expectation, and the main objective of \cite{liu2021orthogonal} is to maintain this low hyperspherical energy throughout training—a factor associated with improved generalization in classification tasks~\cite{liu2018learning,lin2020regularizing}. The findings in \cite{liu2021orthogonal} indicate that such orthogonal transformations possess enough expressivity to produce neural networks with robust generalization for classification. Unlike this line of work, our focus is on adapting text-to-image diffusion models through finetuning, aiming for enhanced control and improved performance on generative tasks. We delineate two fundamental distinctions between OFT and \cite{liu2021orthogonal}. Firstly, whereas \cite{liu2021orthogonal} strives to reduce hyperspherical energy, OFT’s objective is to retain the hyperspherical energy characteristic of the pretrained model, thereby safeguarding its internal structures during finetuning. Notably, in the context of diffusion model adaptation, reducing hyperspherical energy might compromise the essential semantic representations learned in pretraining. Secondly, OFT introduces an explicit mechanism to minimize divergence from the pretrained parameters, leading to a constrained optimization formulation, while \cite{liu2021orthogonal} operates without such regularization. The crux of successful finetuning lies in balancing adaptability with preservation of learned knowledge, and we posit that OFT attains this balance effectively. Our principal contributions can be summarized as follows:

\begin{itemize}[leftmargin=*,nosep]

    \item We introduce Orthogonal Finetuning, an innovative approach designed to enhance the controllability of text-to-image diffusion models during finetuning. Additionally, to bolster model stability, we present a constrained version that restricts the angular difference from the original pretrained parameters.
    \item In contrast with prior finetuning strategies, OFT enables model adaptation while theoretically maintaining the hyperspherical energy. Our experiments indicate that this property is closely linked with preserving the pretrained model's generative semantics.
    \item The effectiveness of OFT is demonstrated across two scenarios: subject-driven image synthesis and generation with explicit controls. Through extensive experimentation, we show that OFT surpasses existing techniques in generation fidelity, training speed, and stability during finetuning. Notably, OFT attains these improvements with higher data and computational efficiency, successfully converging using fewer training samples and epochs.
    \item For tasks involving controllable generation, we propose a novel control consistency metric, which quantitatively measures the model's controllability. This involves estimating the control input from the produced image and comparing it to the original control directive. Our findings show that this metric highlights the robust controllability achieved by OFT.
\end{itemize}

\section{Related Work}

\textbf{Text-to-image diffusion models}. Remarkable advances~\cite{saharia2022photorealistic,ramesh2022hierarchical,rombach2022high,gu2022vector,nichol2022glide} have been achieved in the field of text-to-image synthesis, primarily due to significant strides in diffusion-based generative modeling~\cite{ho2020denoising,song2021score,song2021denoising,dhariwal2021diffusion} and progress in vision-language representation learning~\cite{su2020vlbert,lu2019vilbert,radford2021learning,wang2021simvlm,li2022blip,li2023blip,alayrac2022flamingo,schuhmann2022laion}. Both GLIDE~\cite{nichol2022glide} and Imagen~\cite{saharia2022photorealistic} implement diffusion models operating directly in pixel space. While GLIDE conducts joint training of the text encoder and a diffusion prior utilizing aligned text-image pairs, Imagen instead employs a fixed, pretrained text encoder throughout. In contrast, Stable Diffusion~\cite{rombach2022high} and DALL-E2~\cite{ramesh2022hierarchical} base their diffusion modeling within a latent space framework. The former utilizes VQ-GAN~\cite{esser2021taming} to establish a visual codebook that forms the latent space, whereas the latter leverages CLIP~\cite{radford2021learning} to forge a unified latent embedding space capable of representing both text and images. Beyond diffusion models, research on text-to-image generation has also explored generative adversarial networks~\cite{reed2016generative,zhang2017stackgan,xu2018attngan,li2019controllable} and autoregressive approaches~\cite{ramesh2021zero,ding2021cogview,wu2022nuwa,yuscaling}. OFT functions as a model-agnostic finetuning strategy and is compatible with any text-to-image diffusion model.

\textbf{Subject-driven generation}. To address the challenge of altering subjects, \cite{avrahami2022blended,nichol2022glide} incorporate user-provided masks as supplemental conditions. Context alteration without impacting the subject can also be achieved by leveraging inversion techniques~\cite{choi2021ilvr,dhariwal2021diffusion,ramesh2022hierarchical,gal2022image}. Unlike these mask-dependent approaches, \cite{hertz2022prompt} introduces a method that enables both local and global edits without requiring any masks. However, all the above strategies often struggle to reliably retain fine identity-related attributes of the subject. In contrast, Pivotal Tuning~\cite{roich2022pivotal} addresses identity preservation by finetuning the generator in the vicinity of an initially inverted latent representation and introducing a regularization term targeting identity maintenance. A related technique, \cite{nitzan2022mystyle}, constructs a personalized generative prior for facial imagery by leveraging a dataset specific to an individual. The method presented in \cite{casanova2021instance} generates a range of instance variations but may sacrifice instance-specific nuances. DreamBooth~\cite{ruiz2023dreambooth} advances the text-to-image diffusion model by finetuning it with a dedicated token and a limited set of subject images, employing both reconstruction loss and a class-specific prior preservation loss for identity fidelity. Building on the DreamBooth paradigm, OFT enhances the finetuning procedure by introducing orthogonal transformations, replacing the conventional practice of naïve finetuning at reduced learning rates.

\textbf{Controllable generation}. Image-to-image translation falls within the broader framework of controllable generation, with earlier approaches predominantly utilizing conditional generative adversarial networks~\cite{isola2017image,zhu2017toward,wang2018high,park2019semantic,choi2018stargan}. Recently, diffusion models have also been leveraged for image-to-image translation tasks~\cite{saharia2022palette,voynov2022sketch,wang2022pretraining}. ControlNet~\cite{zhang2023adding} advances this direction by enabling a pretrained diffusion model to be guided through fine-tuning with supplementary control signals, leading to notable improvements in controllability of generation. In parallel, T2I-Adapter~\cite{mou2023t2i} presents a method for fine-tuning pretrained diffusion models to further enhance controllability of the outputs. Adopting the experimental protocols in~\cite{zhang2023adding,mou2023t2i}, we utilize OFT to fine-tune pretrained diffusion models, achieving consistently superior control using reduced amounts of training data and fewer fine-tuned parameters. Importantly, OFT does not incur any extra computational cost during inference.

\textbf{Model finetuning}. The practice of adapting large pretrained models for specific downstream tasks has gained substantial momentum in recent years~\cite{devlin2018bert,brown2020language,he2022masked}. Within this paradigm, adaptation strategies (\emph{e.g.}, \cite{hulora2022,houlsby2019parameter,pfeiffer2020adapterfusion}) have been thoroughly explored in the context of natural language processing. Among these, LoRA~\cite{hulora2022} stands out as particularly related to OFT, operating under the premise that the additive updates to weights during finetuning can be effectively modeled with a low-rank structure. Unlike LoRA, OFT implements a layer-shared orthogonal transformation, which updates the neuron weights multiplicatively. This approach mathematically guarantees the preservation of pair-wise angles among neurons within a layer, thereby contributing to enhanced stability.

\section{Orthogonal Finetuning}

\subsection{Why Does Orthogonal Transformation Make Sense?}

We begin by exploring the motivation for employing orthogonal transformation in the finetuning of text-to-image diffusion models. To clarify this rationale, we break the issue down into two sub-questions: (1) what motivates the adjustment of neuron angles (\emph{i.e.}, direction), and (2) what justifies the use of orthogonal transformation specifically for modifying these angles.

To address the first question, we build upon the empirical findings reported in \cite{liu2018decoupled,chen2020angular}, which indicate that angular differences among features effectively represent the semantic distance between data points. SphereNet~\cite{liu2017deep} further demonstrates that constraining all neurons to the surface of a unit hypersphere during neural network training not only maintains model capacity but can also enhance generalization, suggesting that the orientation of neurons suffices to encapsulate crucial information from the input data. To further clarify the significance of neuron directionality, we perform a controlled experiment depicted in Figure~\ref{fig:toy_ae}. Here, a basic convolutional autoencoder is trained using flower images, where, during training, the feature map is computed through the standard inner product between convolution kernels $\bm{w}$ and sliding window inputs $\bm{x}$, producing outputs $z$. In the evaluation phase, we compare three different methods for generating the feature map: (a) reapplying the inner product as during training, (b) isolating the magnitude information, and (c) extracting the angular information. As illustrated in Figure~\ref{fig:toy_ae}, nearly perfect reconstruction of input images is achieved using angular information alone, whereas magnitude alone fails to preserve semantic content. Notably, the cosine-based angular activation (c) is not utilized during training, as the model is trained exclusively using the inner product. These findings suggest that neuron direction (or angle) primarily encodes the semantic content of the input, indicating that adjusting neuron orientation through finetuning is likely a more effective strategy for manipulating semantic features in images.

Regarding the second question, the most straightforward approach to adjust neuron directions during finetuning is to apply a collective rotation or reflection to all neurons within the same layer, effectively resulting in an orthogonal transformation. While it is possible to consider more adaptable angular transformations, allowing distinct rotations for each neuron, our observations suggest that orthogonal transformation achieves a favorable compromise between adaptability and structure. Additionally, \cite{liu2021orthogonal} provides evidence that orthogonal transformations possess sufficient representational power for neural network training. To substantiate this point, we conducted an experiment assessing the regularization impact of enforcing orthogonality. Utilizing the controllable generation experiment based on the ControlNet~\cite{zhang2023adding} framework, we present outcomes in Figure~\ref{fig:ortho}. These results indicate that the conventional OFT, with orthogonality constraint, exhibits consistent performance and precise controllability upon completing training (at epoch 20). In contrast, the OFT variant lacking this constraint neither produces realistic outputs nor demonstrates any degree of controllability. This experimental comparison underscores the critical role of the orthogonality constraint within OFT.

\subsection{General Framework}

Traditional finetuning approaches generally employ gradient descent with a relatively low learning rate to adjust model parameters, often targeting either the entire network or specific layers. By using a small learning rate, updates are constrained, implicitly limiting the divergence from the original pretrained weights. In effect, classical finetuning pursues an objective that implicitly penalizes large deviations, minimizing $\thickmuskip=2mu \medmuskip=2mu \| \bm{M} - \bm{M}^0\|$, where $\bm{M}$ and $\bm{M}^0$ denote the finetuned and original pretrained model weights, respectively. Although this implicit regularization aligns with intuitive expectations, it may nonetheless allow too much adaptability when applied to large-scale models. LoRA addresses this limitation by incorporating a further restriction: the update to the weights must adhere to a low-rank structure, namely $\thickmuskip=2mu \medmuskip=2mu \text{rank}(\bm{M} - \bm{M}^0)=r'$, where $r'$ is deliberately kept small. In contrast, OFT imposes a different type of constraint—maintaining the similarity of neuron pairs through a hyperspherical energy criterion: $\thickmuskip=2mu \medmuskip=2mu \| \text{HE}(\bm{M})-\text{HE}(\bm{M}^0) \|=0$, where $\text{HE}(\cdot)$ is the hyperspherical energy associated with a weight matrix. For clarity, consider a fully connected layer represented as $\thickmuskip=2mu \medmuskip=2mu \bm{W}=\{\bm{w}_1,\cdots,\bm{w}_n\}\in\mathbb{R}^{d\times n}$, where each column vector $\bm{w}_i\in\mathbb{R}^{d}$ represents a neuron, and $\bm{W}^0$ signifies the pretrained weights. The layer's output $\thickmuskip=2mu \medmuskip=2mu \bm{z}\in\mathbb{R}^n$ is produced via $\thickmuskip=2mu \medmuskip=2mu \bm{z}=\bm{W}^\top\bm{x}$, with $\thickmuskip=2mu \medmuskip=2mu \bm{x}\in\mathbb{R}^d$ as the input. Within this framework, OFT can be understood as enforcing minimization of the hyperspherical energy discrepancy between the updated and pretrained models.

\begin{equation}\label{eq:oft_principle}
\footnotesize
\min_{\bm{W}} \left\| \text{HE}(\bm{W})-\text{HE}(\bm{W}^0)\right\|~~\Leftrightarrow~~\min_{\bm{W}} \bigg{\|} \sum_{i\neq j} \|\hat{\bm{w}}_i-\hat{\bm{w}}_j\|^{-1}-\sum_{i\neq j} \|\hat{\bm{w}}^0_i-\hat{\bm{w}}^0_j\|^{-1}\bigg{\|}
\end{equation}

where $\thickmuskip=2mu \medmuskip=2mu \hat{\bm{w}}_i:=\bm{w}_i/\|\bm{w}_i\|$ represents the $i$-th neuron after normalization, and the concept of hyperspherical energy for a fully connected layer $\bm{W}$ is formalized as $\thickmuskip=2mu \medmuskip=2mu \text{HE}(\bm{W}):= \sum_{i\neq j} \|\hat{\bm{w}}_i-\hat{\bm{w}}_j\|^{-1}$. It is straightforward to see that Eq.~\eqref{eq:oft_principle} reaches a minimum value of zero, which occurs whenever $\bm{W}$ is related to $\bm{W}^0$ by an orthogonal transformation—specifically, $\thickmuskip=2mu \medmuskip=2mu \bm{W}=\bm{R}\bm{W}^0$ for some orthogonal matrix $\thickmuskip=2mu \medmuskip=2mu \bm{R}\in\mathbb{R}^{d\times d}$ (where determinant $1$ and $-1$ correspond to rotation and reflection, respectively). This notion underpins OFT, proposing that neural network adaptation can be accomplished by finetuning via shared orthogonal transformations at the layer level. Concretely, OFT formulates the task as learning the orthogonal matrix $\thickmuskip=2mu \medmuskip=2mu \bm{R}\in\mathbb{R}^{d \times d}$ given a pretrained fully connected layer $\thickmuskip=2mu \medmuskip=2mu\bm{W}^0\in\mathbb{R}^{d\times n}$, thereby modifying the forward computation from $\thickmuskip=2mu \medmuskip=2mu\bm{z}=(\bm{W}^0)^\top\bm{x}$ to

\begin{equation}\label{eq:oft_general}
    \footnotesize
    \bm{z}=\bm{W}^\top\bm{x}=(\bm{R}\cdot\bm{W}^0)^\top\bm{x},~~~\text{s.t.}~\bm{R}^\top\bm{R}=\bm{R}\bm{R}^\top=\bm{I}
\end{equation}

Here, $\bm{W}$ represents the weight matrix fine-tuned via OFT, while $\bm{I}$ stands for the identity matrix. Figure~\ref{fig:block} provides an illustration of OFT. In parallel to the approach of zero initialization used in LoRA, it is important for OFT to begin fine-tuning from the original pretrained parameter values, denoted as $\bm{W}^0$. This requirement is satisfied by initializing the orthogonal matrix $\bm{R}$ as an identity matrix, thereby ensuring that the initial state of the fine-tuned model matches the pretrained configuration. Maintaining the orthogonality of $\bm{R}$ can be addressed through differentiable orthogonalization methods described in \cite{liu2021orthogonal,lezcano2019cheap}. The following sections will elaborate on ensuring orthogonality in a manner that is computationally efficient.

\subsection{Efficient Orthogonal Parameterization}\label{sect:eff_ortho}

Conventional orthogonalization techniques like the Gram-Schmidt process, although differentiable, typically incur substantial computational costs~\cite{liu2021orthogonal}. To achieve greater computational efficiency, we leverage Cayley parameterization for constructing orthogonal matrices. In particular, our approach defines the orthogonal matrix as $\thickmuskip=2mu \medmuskip=2mu \bm{R}=(\bm{I}+\bm{Q})(I-\bm{Q})^{-1}$, with $\bm{Q}$ denoting a skew-symmetric matrix where $\thickmuskip=2mu \medmuskip=2mu \bm{Q}=-\bm{Q}^\top$. While this method enhances efficiency, it introduces a minor drawback: the Cayley parameterization inherently restricts us to orthogonal matrices with a determinant of $1$, confining $\bm{R}$ to the special orthogonal group. Nevertheless, our empirical results indicate that this restriction has negligible impact on practical performance. Moreover, despite using the Cayley transform for $\bm{R}$, the parameterization can become highly inefficient as dimensionality $d$ increases. To mitigate this issue, we propose representing $\bm{R}$ as a block-diagonal matrix divided into $r$ blocks, resulting in the following structure:

\begin{equation}
    \footnotesize
    \bm{R}=\text{diag}(\bm{R}_1,\bm{R}_2,\cdots,\bm{R}_r)=\begin{bmatrix}
    \bm{R}_1\in \text{O}(\frac{d}{r}) &  &\\
    &\ddots & \\
    & & \bm{R}_r\in \text{O}(\frac{d}{r})
    \end{bmatrix}\in \text{O}(d)
\end{equation}

where $\text{O}(d)$ represents the orthogonal group of degree $d$, and $\thickmuskip=2mu \medmuskip=2mu \bm{R}\in\mathbb{R}^{d\times d}$ as well as $\thickmuskip=2mu \medmuskip=2mu \bm{R}_i\in\mathbb{R}^{d/r\times d/r},\forall i$ are orthogonal matrices. In the particular case where $\thickmuskip=2mu \medmuskip=2mu r=1$, the block-diagonal orthogonal matrix reduces to a general unconstrained orthogonal matrix. An orthogonal matrix of dimension $\thickmuskip=2mu \medmuskip=2mu d\times d$ has $\thickmuskip=2mu \medmuskip=2mu d(d-1)/2$ independent parameters, resulting in an overall complexity of $\mathcal{O}(d^2)$. Conversely, for an $r$-block diagonal orthogonal matrix, the parameter count is given by $\thickmuskip=2mu \medmuskip=2mu d(d/r-1)/2$, leading to a complexity of $\thickmuskip=2mu \medmuskip=2mu \mathcal{O}(d^2/r)$. One may further reduce the parameter count by sharing the block matrices, \emph{i.e.}, letting $\thickmuskip=2mu \medmuskip=2mu\bm{R}_i=\bm{R}_j$ for all $i\neq j$, which lowers the complexity to $\mathcal{O}(d^2/r^2)$. Importantly, despite these modifications designed for improved parameter efficiency, the overall transformation $\bm{R}$ remains an orthogonal matrix, thus preserving hyperspherical energy without compromise.

We analyze the parameter efficiency of OFT relative to LoRA. In LoRA, given a low-rank setting with rank $r'$, the count of trainable parameters can be expressed as $\thickmuskip=2mu \medmuskip=2mu r'(d+n)$. When both $r$ and $r'$ are proportional to the neuron dimension $d$ (\emph{e.g.}, $\thickmuskip=2mu \medmuskip=2mu r=r'=\alpha d$ for some fixed $\alpha$ in the interval $0<\alpha\leq 1$), the parameter scaling for LoRA is $\mathcal{O}(d^2+dn)$, whereas OFT exhibits a parameter scaling of only $\mathcal{O}(d)$. To concretely compare these complexities, consider a case with a weight matrix of dimensions $\thickmuskip=2mu \medmuskip=2mu 128\times 128$: under LoRA with $\thickmuskip=2mu \medmuskip=2mu r'=8$, there are $2,048$ parameters that need to be trained; in contrast, OFT—using the same value of $\thickmuskip=2mu \medmuskip=2mu r=8$ and without block sharing—requires only $960$ trainable parameters.

\subsection{Constrained Orthogonal Finetuning}

The original OFT approach can be made more restrictive by ensuring that the finetuned model remains closely aligned with the pretrained model. To accomplish this, COFT adopts the following method during the forward pass:

\begin{equation}\label{eq:coft_general}
    \footnotesize
    \bm{z}=\bm{W}^\top\bm{x}=(\bm{R}\cdot\bm{W}^0)^\top\bm{x},~~~\text{s.t.}~\bm{R}^\top\bm{R}=\bm{R}\bm{R}^\top=\bm{I},~\norm{\bm{R}-\bm{I}}\leq \epsilon
\end{equation}

which imposes both an orthogonality constraint and a requirement that deviations from the identity matrix are limited by $\epsilon$. The orthogonality condition can be satisfied using the Cayley parameterization as discussed in Section~\ref{sect:eff_ortho}. Nonetheless, integrating the $\epsilon$-deviation condition within the Cayley-parametrized orthogonal matrix is not straightforward. To develop a better understanding of the Cayley transform, we leverage the Neumann series for approximating $\thickmuskip=2mu \medmuskip=2mu \bm{R}=(\bm{I}+\bm{Q})(I-\bm{Q})^{-1}$, which yields $\thickmuskip=2mu \medmuskip=2mu \bm{R}\approx\bm{I}+2\bm{Q}+\mathcal{O}(\bm{Q}^2)$, assuming that the Neumann {\looseness=-1 \parfillskip=0pt\par}

As a consequence, the requirement $\thickmuskip=2mu \medmuskip=2mu\|\bm{R}-\bm{I}\|\leq\epsilon$ can be incorporated within the Cayley transform framework. This condition is equivalently translated to $\thickmuskip=2mu \medmuskip=2mu \|\bm{Q}-\bm{0}\|\leq\epsilon'$, where $\bm{0}$ is the matrix consisting entirely of zeros and $\epsilon'$ is a distinct error hyperparameter from $\epsilon$. This reformulated constraint on $\bm{Q}$ is readily manageable through projected gradient descent methods. To ensure that the orthogonal matrix $\bm{R}$ is initialized as the identity, we simply begin with $\bm{Q}$ as the zero matrix. The COFT approach unifies two direct constraints: maintaining a small hyperspherical energy discrepancy and explicitly restricting the deviation from the original pretrained model. While the latter constraint is often tacitly employed in prior finetuning strategies, COFT formalizes it within its objective. Even though OFT has shown remarkable efficacy, our empirical findings suggest that COFT offers improved stability during finetuning, largely attributable to the explicit enforcement of this deviation bound. Figure~\ref{fig:eps_coft} illustrates how the parameter $\epsilon$ modulates COFT's performance. Specifically, as $\epsilon$ increases, the behavior of the COFT-finetuned model more closely matches that produced by OFT, whereas decreasing $\epsilon$ results in outputs that are progressively more akin to those generated by the original pretrained text-to-image diffusion model.

\subsection{Re-scaled Orthogonal Finetuning}

We introduce a straightforward modification to the standard OFT approach by incorporating a learnable scaling factor for each neuron. This addition is inspired by the observation that scaling the output of neurons does not alter the hyperspherical energy, as the normalization process removes the effect of magnitude. Concretely, our proposed forward computation is given by $\thickmuskip=2mu \medmuskip=2mu\bm{z}=(\bm{R}\bm{W}^0\bm{D})^\top\bm{x}$\footnote{\textbf{Errata}: The NeurIPS camera-ready version incorrectly reported the forward pass for the re-scaled OFT as $\thickmuskip=2mu \medmuskip=2mu\bm{z}=(\bm{D}\bm{R}\bm{W}^0)^\top\bm{x}$. However, the original implementation used the correct formulation, and the results shown in Appendix~\ref{app:rescaled} remain valid.}, where $\thickmuskip=2mu \medmuskip=2mu \bm{D}=\text{diag}(s_1,\cdots,s_n)\in\mathbb{R}^{n\times n}$ is a diagonal matrix composed of learnable entries $s_1,\cdots,s_n$ (all strictly positive). Unlike the original OFT, whose forward pass in Eq.~\eqref{eq:oft_general} restricts learning to $\bm{R}$, our variant jointly learns both $\bm{D}$ and the orthogonal matrix $\bm{R}$. This re-scaled OFT offers enhanced modeling flexibility at the cost of only a trivial increase in parameter count. While our experiments utilize the vanilla OFT to isolate the impact of orthogonal transformations, we observe that the re-scaled OFT typically outperforms the original variant (refer to Appendix~\ref{app:rescaled}).

\section{Intriguing Insights and Discussions}

\textbf{OFT is compatible with various neural network architectures}. In theory, OFT can be utilized with any neural network model. For example, in Transformers, LoRA is generally utilized to finetune the attention weight matrices~\cite{hulora2022}. In order to conduct a fair comparison with LoRA, our experimental setup restricts OFT to finetune only the attention weights. Beyond standard fully connected layers, OFT is particularly amenable to finetuning convolutional layers as well, owing to the unique block-diagonal arrangement of $\bm{R}$, which has meaningful interpretations in the context of convolutions—a feature not paralleled by LoRA. If the quantity of blocks equals the number of input channels, each block operates on a distinct neuron channel, mirroring the mechanism of depth-wise convolution kernel learning~\cite{chollet2017xception}. Conversely, if all blocks of $\bm{R}$ are tied together, the transformation applies the same orthogonal matrix across all channels. Additional explorations involving OFT for convolutional layer finetuning are provided in Appendix~\ref{app:convolution}.

\textbf{Connection to LoRA}. LoRA achieves stability in the pretrained weight matrix by introducing a low-rank update, thereby limiting significant changes to the existing information. In comparison, OFT maintains the integrity of pretrained weights by enforcing the update transformation to be orthogonal (and thus full-rank), ensuring that pretraining information is not compromised by the adaptation. The forward computation in OFT can be reformulated as $\thickmuskip=2mu \medmuskip=2mu \bm{z}=(\bm{R}\bm{W}^0)^\top\bm{x}=(\bm{W}^0+(\bm{R}-\bm{I})\bm{W}^0)^\top\bm{x}$, where the term $\thickmuskip=2mu \medmuskip=2mu (\bm{R}-\bm{I})\bm{W}^0$ plays a similar role to LoRA’s low-rank weight update. Since $\bm{W}^0$ is generally of full rank, OFT also implements a low-rank adaptation when $\thickmuskip=2mu \medmuskip=2mu \bm{R}-\bm{I}$ is low-rank. Both methods feature hyperparameters to control parameter efficiency: LoRA’s rank parameter $r'$ and OFT’s diagonal block parameter $r$ allow for reducing the set of trainable weights. Importantly, LoRA and OFT reflect fundamentally different principles of parameter efficiency; whereas LoRA utilizes low-rank approximations to minimize trainable parameters, OFT leverages the sparsity of block-diagonal orthogonal matrices to achieve the same goal.

\textbf{Why OFT converges faster}. As illustrated in Figure~\ref{fig:toy_ae}, adjusting neuron directions provides the most efficient means for semantic modification, aligning precisely with the main mechanism behind OFT. Additionally, OFT can be interpreted as an approach that finetunes neurons on a smooth hyperspherical manifold, thereby facilitating a more favorable optimization landscape. Empirical support for this is also presented in \cite{liu2021orthogonal}.

\textbf{Reason for not pursuing hyperspherical energy minimization}. Unlike the approach taken in \cite{liu2021orthogonal}, we refrain from minimizing hyperspherical energy. Classification tasks benefit from non-redundant neurons; however, achieving minimum hyperspherical energy leads to neurons being distributed evenly across the hypersphere. For finetuning purposes, this objective is not appropriate, as it can potentially disrupt the valuable information acquired during pretraining.

\textbf{Balancing flexibility and regularity during finetuning}. Our findings reveal a fundamental tension between flexibility and regularity. Although standard finetuning offers maximal adaptability, it often compromises stability and can readily lead to model degeneration. Interestingly, the straightforward approach of OFT manages to strike an effective equilibrium, maintaining pairwise neuron angles to balance adaptability with structural preservation. The block-diagonal parameterization serves as a more rigorous regularization mechanism for the orthogonal matrix, further constraining the model.

\textbf{No extra inference cost}. In contrast to ControlNet, the OFT framework does not incur any additional computational burden during inference for the fine-tuned model. At inference time, it is sufficient to incorporate the learned orthogonal matrix $\bm{R}$ by multiplying it with the original pretrained weights $\bm{W}^0$, resulting in a new, equivalent weight matrix $\thickmuskip=2mu \medmuskip=2mu \bm{W}=\bm{R}\bm{W}^0$. As a result, the model’s inference speed remains unchanged from that of the pretrained version.

\section{Experiments and Results}

\textbf{General settings}. Throughout our experiments, we utilize Stable Diffusion v1.5~\cite{rombach2022high} as the foundation for the pretrained text-to-image generation tasks. To maintain fairness across methods, image samples are selected randomly from the outputs of each technique. The protocol for subject-driven image creation is adapted from DreamBooth~\cite{ruiz2023dreambooth}, while procedures for controlled generation are based on ControlNet~\cite{zhang2023adding} and T2I-Adapter~\cite{mou2023t2i}. For direct comparison with LoRA, OFT and COFT modifications are confined to the identical layers utilized by LoRA. Additional experimental results and implementation specifics are included in Appendix~\ref{app:settings}.

\subsection{Subject-driven Generation}

\textbf{Settings}. As baselines, we employ both DreamBooth~\cite{ruiz2023dreambooth} and LoRA~\cite{hulora2022}. Each method utilizes the same loss function prescribed by DreamBooth. For both DreamBooth and LoRA, we primarily adhere to the configurations recommended by their respective original works, opting for their optimal hyperparameters. Additional experimental results can be found in Appendix~\ref{app:settings},\ref{app:coft_oft},\ref{app:more_qualitative},\ref{app:failure}.

\textbf{Finetuning stability and convergence}. We begin by assessing the stability and convergence rate during finetuning for DreamBooth, LoRA, OFT, and COFT. The findings, as illustrated in Figure~\ref{fig:teaser} and Figure~\ref{fig:db_convergence}, indicate that COFT and OFT both allow the diffusion model to be finetuned with considerable stability. Notably, after 400 iterations, both DreamBooth and OFT-based approaches demonstrate effective control, whereas LoRA is unable to maintain the subject’s identity. By the 2000-iteration mark, DreamBooth exhibits mode collapse—producing degenerate images—and LoRA fails to render the yellow shirt, frequently producing yellow fur instead. In comparison, OFT and COFT consistently maintain stable and accurate control over the outputs throughout this process. These observations affirm that our OFT framework enables subject-driven generation with rapid yet reliable convergence. It is worth emphasizing that the performance of OFT and COFT is largely insensitive to the choice of the iteration count, which removes the necessity of precise tuning for different subjects. As a result, one can simply select a sufficiently high number of finetuning steps for both OFT and COFT without concern. For COFT, given a suitable value of $\epsilon$, both the learning rate and number of iterations can be chosen without substantial effort.

\textbf{Quantitative comparison}. In line with the methodology presented in \cite{ruiz2023dreambooth}, we quantitatively assess subject fidelity (using DINO~\cite{caron2021emerging} and CLIP-I~\cite{radford2021learning}), fidelity to text prompts (using CLIP-T~\cite{radford2021learning}), as well as sample diversity (with LPIPS~\cite{zhang2018unreasonable}). The CLIP-I metric measures the mean pairwise cosine similarity of CLIP embeddings between generated samples and their corresponding real images. Analogously, the DINO score is calculated in the same fashion, except it leverages ViT S/16 DINO embeddings rather than CLIP. For evaluating text prompt fidelity, CLIP-T computes the average cosine similarity between the CLIP embeddings of the generated images and their associated textual prompts. Additionally, we measure sample diversity as the average LPIPS cosine similarity across generated images pertaining to the same subject and text prompt. As presented in Table~\ref{table:dreambooth_final}, COFT and OFT both significantly surpass DreamBooth and LoRA in terms of DINO and CLIP-I metrics, while also delivering slightly superior or comparable results on the text prompt fidelity and diversity metrics. Each method is finetuned multiple times, specifically using 30 distinct random seeds per pretrained model, to mitigate variance from stochastic effects. Collectively, the findings highlight that our OFT framework delivers superior convergence behavior, greater training stability, and consistently outperforms alternatives in final metric outcomes.

\textbf{Qualitative comparison}. To provide a clearer insight into the advantages offered by OFT, we present a set of randomly selected samples for subject-driven generation, as illustrated in Figure~\ref{fig:dreambooth_final}. For consistency and fairness, all images are synthesized using the same fine-tuned model architecture for each method, without any individualized optimization of text prompts for final outputs. For every method under comparison, we choose the checkpoint yielding the highest validation CLIP scores. Upon inspecting the examples in Figure~\ref{fig:dreambooth_final}, it becomes apparent that both OFT and COFT excel at retaining semantic subject identity, in contrast to LoRA, which frequently struggles in this regard (\emph{e.g.}, the subject’s features are entirely lost by LoRA in the bowl case). Additionally, OFT and COFT demonstrate substantially better alignment with the textual prompts, whereas DreamBooth—despite successfully preserving subject characteristics—often does not adhere closely to prompt content (\emph{e.g.}, see the first row for the bowl example). Collectively, these qualitative results indicate that the OFT approach offers superior prompt controllability alongside strong subject preservation. Furthermore, OFT demonstrates robustness to the choice of iteration numbers, performing effectively even as iterations increase, in contrast to DreamBooth and LoRA, which lack such resilience. Additional qualitative illustrations can be found in Appendix~\ref{app:more_qualitative}. Our human assessment results, detailed in Appendix~\ref{app:human}, provide further evidence of OFT’s effectiveness.

\subsection{Controllable Generation}

\textbf{Settings}. Our experimental baselines include ControlNet~\cite{zhang2023adding}, T2I-Adapter~\cite{mou2023t2i}, and LoRA~\cite{hulora2022}. In the primary experiments, we address three demanding controllable generation scenarios: converting Canny edge maps to images~(C2I) using the COCO dataset~\cite{lin2014microsoft}, generating images from segmentation maps~(S2I) with the ADE20K dataset~\cite{zhou2017scene}, and transforming landmark annotations to facial images~(L2F) with CelebA-HQ~\cite{karras2018progressive,xia2021tedigan}. For each task, the respective methods are employed to finetune Stable Diffusion~(SD) v1.5 on the corresponding datasets for a total of 20 epochs. Additional findings are presented in Appendix~\ref{app:more_qualitative},\ref{app:more_control_task},\ref{app:failure}.

\textbf{Convergence}. We analyze the convergence characteristics of ControlNet, T2I-Adapter, LoRA, and COFT within the context of the L2F task, employing both quantitative metrics and qualitative assessments. For the quantitative evaluation, we measure the average $\ell_2$ norm between predicted facial landmarks and their corresponding control landmarks. Figure~\ref{fig:face_convergence} presents the trend of facial landmark error across various epochs for each model during finetuning. Notably, both COFT and OFT demonstrate a markedly accelerated convergence rate. For instance, while LoRA requires 20 epochs to match the performance level that OFT achieves in just 8 epochs. It is important to highlight that both OFT and COFT possess a comparable number of trainable parameters to LoRA—substantially fewer than those used by ControlNet—yet they exhibit a considerably higher convergence efficiency than the competing approaches. The expedited convergence of OFT is further supported by the evidence in Figure~\ref{fig:teaser}. The rightmost example illustrates that OFT attains robust convergence using only 5\% of the ADE20K dataset, underscoring its data efficiency relative to ControlNet and LoRA. In terms of qualitative comparison, our analysis centers on OFT, COFT, and ControlNet, since the latter attains landmark errors most similar to those of our methods. As depicted in Figure~\ref{fig:face_convergence_qual}, both OFT and COFT consistently and smoothly align the synthesized facial poses with the control landmarks over time. Conversely, ControlNet exhibits an abrupt improvement in controllability beginning at the 8-th epoch, echoing the convergence behavior reported in \cite{zhang2023adding}. The pronounced stability and predictability of OFT's convergence dynamics facilitate more reliable and manageable training in practical applications, making the process of selecting effective hyperparameters for OFT substantially more straightforward.

\textbf{Quantitative comparison}. To assess controllable generation, we propose a control consistency metric. This metric involves extracting the control signal from the generated output and measuring its alignment with the initial control input. Specifically, for the C2I task, we use IoU and F1 score as our evaluation criteria. In the S2I task, metrics such as mean IoU, mean accuracy, and overall accuracy are employed. For the L2F task, we evaluate performance using the mean $\ell_2$ distance between the predicted and original control landmarks. Further specifics on these consistency metrics can be found in Appendix~\ref{app:settings}. Each method in comparison is run with optimized hyperparameter configurations. According to the results presented in Table~\ref{table:control_gen}, both OFT and COFT consistently outperform alternative methods in terms of control fidelity and precision. In contrast, adapter-based strategies (including T2I-Adapter and ControlNet) demonstrate slower convergence and inferior final outcomes. LoRA achieves stronger performance on the S2I task compared to ControlNet, yet it lags behind on C2I and L2F tasks. Generally, we observe that LoRA's maximum achievable performance remains limited, even with extensive hyperparameter optimization. Conversely, OFT’s performance continues to improve with an increased parameter count and has yet to plateau. We highlight that our proposed quantitative assessment for controllable generation constitutes a significant contribution, as it enables precise evaluation of control quality for finetuned models—bounded only by the capabilities of standard segmentation or detection models used for measurement.

\textbf{Qualitative comparison}. We further provide a qualitative assessment of OFT and COFT in relation to leading contemporary approaches such as ControlNet, T2I-Adapter, and LoRA. As depicted by the randomly sampled outputs in Figure~\ref{fig:control_final}, OFT and COFT consistently generate images with superior fidelity and realism, alongside demonstrating precise controllability. Within the S2I task, LoRA fails to adhere to the structural constraints imposed by the input segmentation map, whereas ControlNet, OFT, and COFT are capable of reliably controlling the synthesized outputs. Remarkably, compared to ControlNet, OFT and COFT deliver images with enhanced detail and plausible geometry, all while utilizing significantly fewer model parameters. In the context of the C2I task, OFT and COFT successfully produce images with strong semantic alignment from coarse Canny edge maps, a capability in which T2I-Adapter and LoRA underperform. Regarding the L2F task, our framework yields the most precise facial pose control, even for challenging head orientations. Across all three examined control tasks, both OFT and COFT consistently surpass existing methods in qualitative image quality, underscoring the potency of the OFT architecture for controlled image generation. To ensure a thorough qualitative evaluation, supplementary results across all control tasks are provided in Appendix~\ref{app:control_exp}, while Appendix~\ref{app:more_control_task} illustrates the adaptability of OFT to further control scenarios, such as dense pose to human body, sketch to image, and depth to image conversions.

\section{Concluding Remarks and Open Problems}

Building upon the insight that the angular relationships between neurons play a key role in shaping visual semantics, we introduce orthogonal finetuning—a straightforward yet powerful technique for finetuning text-to-image diffusion models. Our approach addresses two primary tasks: subject-driven generation and controllable generation. Relative to current approaches, OFT achieves superior control and greater stability during finetuning, all while utilizing a reduced set of parameters. Notably, OFT incurs no extra computational burden during inference, resulting in a model that is both efficient and practical for deployment.

OFT presents several noteworthy open questions. To begin with, OFT ensures orthogonality through the use of Cayley parametrization, which necessitates the computation of a matrix inverse and can pose scalability issues. Although block diagonal parametrization partially alleviates this restriction, developing differentiable methods to further accelerate the matrix inversion remains unresolved. Furthermore, OFT offers intriguing prospects for compositionality: orthogonal matrices generated through independent OFT finetuning processes can be composed by multiplication, resulting in another orthogonal matrix. However, it is not yet clear whether the composed orthogonal matrices adequately retain knowledge from all participating downstream tasks, presenting a compelling avenue for further research. Lastly, the parameter efficiency achieved by OFT primarily relies on the adoption of block diagonal structures, which can introduce systematic biases and limit expressive capacity. Identifying strategies to enhance parameter efficiency without incurring such biases is a significant and unresolved challenge.

\end{document}